\documentclass[11pt,a4paper]{article}

\usepackage[margin=2.2cm]{geometry}
\usepackage[T1]{fontenc}
\usepackage[utf8]{inputenc}
\usepackage{lmodern}
\usepackage{microtype}
\usepackage{booktabs}
\usepackage{tabularx}
\usepackage{array}
\usepackage{amsmath}
\usepackage{xcolor}
\usepackage[hidelinks]{hyperref}
\hypersetup{pdftitle={Supporting Information: Beyond Candidate Generation: A Validation-First Protocol for AI-Assisted Transition-State Discovery}, pdfauthor={Haofan Li}}
\usepackage{caption}
\captionsetup{font=small,labelfont=bf}
\setlength{\parindent}{1.1em}
\setlength{\parskip}{0.25em}
\renewcommand{\arraystretch}{1.12}

\title{Supporting Information\\[4pt]
\large Beyond Candidate Generation: A Validation-First Protocol for AI-Assisted Transition-State Discovery}
\author{Haofan Li}
\date{}

\renewcommand{\thetable}{S\arabic{table}}
\begin{document}
\maketitle

\section*{S1. Protocol checklists}

\subsection*{S1.1 P0: chemical preflight}
\begin{itemize}
\item Reactive pairs: which bonds form and which break, with their actual distances in the starting geometry.
\item Charge, multiplicity, and every accessible $d^n$ spin state, including broken-symmetry solutions.
\item Connectivity: build a bond graph with a distance criterion; check the intended bonds and any unplanned bond.
\item Conformers: pre-fold flexible ligands/chains; retain several low-energy conformers.
\item Reference state: reactant or encounter complex and product, at the same level as the TS.
\item Candidate filter: hard overlaps ($<1.00$~\AA{}), clash ratios ($<0.75\,\Sigma r_{\mathrm{cov}}$), heavy-atom RMSD, connectivity changes.
\end{itemize}

\subsection*{S1.2 P0.5: learned-potential pre-screen (optional)}
The pre-screen sits between P0 and P1 and is the only stage that runs on a
machine-learned potential (MLIP) instead of on DFT. It is optional, and no gate that
decides whether a candidate is a transition state depends on it.

\begin{itemize}
\item Declare the reaction coordinate before screening: the forming and breaking pairs, the target bond lengths, the charge, and the multiplicity. The screen changes the geometry; it does not change the electronic-state assignment.
\item Use the potential only inside its applicability domain. For broken-symmetry and open-shell transition-metal systems the verdict describes the learned surface, not the electronic state, and is recorded as a PES mismatch rather than as a basin.
\item Displace the candidate along the declared coordinate in both directions and relax on the learned surface (the reference implementation defaults to $0.35$~\AA{} per side).
\item Read the pair of endpoints. Same minimum on both sides: basin-like, discard before any DFT is spent. The two minima differ by exactly the declared forming and breaking bonds: barrier-like, pass the relaxed geometry to P1. Any other difference: the reactive pairs are wrong, and S1.1 is repeated before the screen is run again.
\item PES mismatch: if the candidate lies more than 60~kcal~mol$^{-1}$ above both endpoints, the learned surface does not support this reaction. Do not read the verdict as a basin; skip the screen and go to P1 (rule 11 of Table~1 of the main text).
\item Geometry repair: remove hard overlaps (0.83--0.93~\AA{} in the cases reported here) and pull the declared bonds to their target lengths while the coordination sphere is held. Check that the coordination bonds are unchanged to within 0.05~\AA{} before the geometry is used.
\item Do not budget the repair as saved optimization cycles: ORCA projects the constraints onto the input geometry before the first cycle (measured, Ni--C$^*$ 3.000 $\rightarrow$ 2.150~\AA{}), and the cycles that follow relax the ligand conformers.
\item Record the model identity and version, the device, the declared pairs, and the verdict with the candidate, so that the pre-screen can be repeated or audited.
\end{itemize}

\subsection*{S1.3 P1: barrier-region scan}
\begin{itemize}
\item Use typically 4--8 points along the reaction coordinate (more when the channel is shallow, has a double bump, or the first scan is inconclusive; the case study in the main text used 16), not a blind OptTS.
\item Each point is an independent constrained optimization, starting from the previous converged geometry.
\item Only converged points enter the profile; the highest converged point is the TS guess.
\item Check the profile for monotonicity, shallow double bumps, and flat regions.
\end{itemize}

\subsection*{S1.4 P2: constrained reference and Hessian}
\begin{itemize}
\item Run a constrained \texttt{Opt + Freq} at the chosen distance.
\item The lowest modes must involve the declared reactive pair; otherwise do not run OptTS.
\item Keep the Hessian and the reference geometry; both are reused downstream.
\end{itemize}

\subsection*{S1.5 P3: OptTS}
\begin{itemize}
\item Read the P2 Hessian (\texttt{InHess Read} in ORCA; \texttt{readfc} in Gaussian).
\item If displacing, scale by the target bond-length change ($\sim$0.05~\AA{}), not the maximum atomic displacement.
\item Apply the stop rules collected in Table~1 of the main text.
\end{itemize}

\subsection*{S1.6 P4: validation}
\begin{itemize}
\item Exactly one significant imaginary frequency.
\item The imaginary mode moves the declared reactive pair; localization is not below 15\%.
\item Key bond lengths match the intended step.
\item IRC or downhill relaxation reaches the intended reactant and product.
\item $\langle S^2\rangle$ and Mulliken spin are consistent with the geometry.
\item Optimization/frequency/IRC at the same level; $\Delta E$ for minimum-vs-saddle; one solvation convention.
\end{itemize}

\section*{S2. Input templates}

All templates use the level of theory of Section S3 unless stated. \texttt{\%maxcore 1100} is a per-core limit (1100~MB), so the twelve-core examples need about 13~GB of scratch space; ORCA stops with an error when a per-core limit cannot be satisfied.

\subsection*{S2.1 ORCA: per-point constrained optimization}
\begin{verbatim}
! Opt TPSS D3BJ def2-SVP RIJCOSX
%pal nprocs 12 end
%maxcore 1100
%cpcm
  smd true
  SMDsolvent "Toluene"
end
%geom
  MaxIter 220
  Constraints
    { B 20 56 1.750 C }
  end
end
* xyzfile 0 2 start.xyz
\end{verbatim}

\subsection*{S2.2 ORCA: constrained reference + frequency}
\begin{verbatim}
! Opt Freq TPSS D3BJ def2-SVP RIJCOSX
%pal nprocs 12 end
%maxcore 1100
%cpcm
  smd true
  SMDsolvent "Toluene"
end
%geom
  MaxIter 220
  Constraints
    { B 20 56 1.750 C }
  end
end
* xyzfile 0 2 start.xyz
\end{verbatim}

\subsection*{S2.3 ORCA: OptTS reading the reference Hessian}
\begin{verbatim}
! OptTS Freq TPSS D3BJ def2-SVP RIJCOSX
%pal nprocs 12 end
%maxcore 1100
%cpcm
  smd true
  SMDsolvent "Toluene"
end
%geom
  MaxIter 160
  InHess Read
  InHessName "L1-constOptFreq.hess"
end
* xyzfile 0 2 L1-constOptFreq.xyz
\end{verbatim}

\subsection*{S2.4 ORCA: IRC}
\begin{verbatim}
! IRC TPSS D3BJ def2-SVP RIJCOSX
%pal nprocs 12 end
%maxcore 1100
%irc
  MaxIter 20
  InitHess Read
  Hess_Filename "L2-TSopt.hess"
  Direction both
  PrintLevel 2
end
* xyzfile 0 2 L2-TSopt.xyz
\end{verbatim}

\subsection*{S2.5 Gaussian: constrained reference + frequency}
\begin{verbatim}
%mem=48GB
%nprocshared=24
%chk=01_const_optfreq.chk
#p UB3LYP/6-31G(d) opt=(calcfc,modredundant,maxcycles=220) freq
   em=gd3bj scrf=(smd,solvent=toluene) int=ultrafine

TS step 1: constrained Opt+Freq

0 2
<coordinates>

B 1 2 1.500 F
B 3 4 2.100 F

\end{verbatim}

The level of theory in the Gaussian templates is an example; substitute the level used in the production runs (the case studies used TPSS-D3BJ/def2-SVP with SMD(toluene), see Section S3).

\subsection*{S2.6 Gaussian: OptTS reading the reference checkpoint}
\begin{verbatim}
%mem=48GB
%nprocshared=24
%oldchk=01_const_optfreq.chk
%chk=02_ts_optfreq.chk
#p UB3LYP/6-31G(d) opt=(ts,readfc,noeigen,maxcycles=220) freq
   em=gd3bj scrf=(smd,solvent=toluene) int=ultrafine
   geom=check guess=read

TS step 2: OptTS+Freq

0 2

\end{verbatim}

\subsection*{S2.7 Slurm template}
\begin{verbatim}
#!/bin/bash
#SBATCH -J ts
#SBATCH -N 1
#SBATCH -n 24
#SBATCH -p <partition>
#SBATCH -o slurm-%j.out
#SBATCH --time=48:00:00

export g16root='<g16root>'
source "${g16root}/g16/bsd/g16.profile"
export GAUSS_SCRDIR=/tmp

date
bash run_chain.sh
echo "EXIT=$?"
date
\end{verbatim}

\subsection*{S2.8 Learned-potential pre-screen and geometry repair}
Both tools belong to the same package, run without a quantum-chemistry engine,
and are called on an \texttt{xyz} candidate. The reactive pairs are given
explicitly and the model is named on the command line, so that a reported
verdict can be reproduced:
\begin{verbatim}
# basin-like / barrier-like verdict (pairs are 1-based)
python scripts/ts_mlip_screen.py \
    --cand candidate.xyz \
    --forming 21,72 --breaking 1,21 \
    --charge 0 --spin 2 \
    --model orb_v3_conservative_omol --device cuda

# geometry repair: i,j,d are 0-based, d in A
python scripts/mlip_preopt.py \
    --in raw.xyz --out repaired.xyz \
    --constraint 20,56,1.75 \
    --freeze-shell 1
\end{verbatim}
\texttt{ts\_mlip\_screen.py} reports one of \texttt{BARRIER-LIKE},
\texttt{BASIN}, \texttt{WRONG-CHANNEL}, or \texttt{PES-MISMATCH} per candidate;
with \verb|--bs| (broken-symmetry singlet) it restricts itself to the geometry
and issues no connectivity verdict. Indices are 1-based there and 0-based in
\texttt{mlip\_preopt.py}, which follows the convention of the ORCA
\texttt{\%geom Constraints} block. A run returns either a repaired geometry or a
pair of relaxation endpoints; neither is a transition-state judgement, and the
endpoints are read as described in S1.2.

\section*{S3. Case-study details}

The systems in the main-text case studies are 84--98 atom Ni/Fe complexes with
organic ligands, treated with TPSS-D3BJ/def2-SVP (RIJCOSX, SMD) for geometry
and frequency calculations, and with higher-level single points where noted.
The numbers quoted below are the observations discussed in the main text; the
case numbers are those of Table~3 there. RIJCOSX is retained because it appears
in the production inputs; for a non-hybrid functional the COSX approximation
acts on no Hartree--Fock exchange.


\section*{S4. AI candidate generation: practical notes}

\subsection*{S4.1 UniTS-Gen}
The UniTS code package (which implements the UniTS-Gen model) takes \texttt{reactive\_atom\_idx} as 0-based, whereas many post-processing
scripts use 1-based atom indices. Convert explicitly. The intended use is
multi-sampling: generate ten candidates, prefilter, and send only a few to DFT.
Single-shot correctness is moderate; the package's published benchmark reports that
multi-sampling recovers all five previously failed cases. In an internal test
on an 84-atom Ni broken-symmetry system, four candidates passed the geometric
filter but had 3.6--4.4~\AA{} heavy-atom RMSD to the reference; the model was out
of distribution for that system class.

\subsection*{S4.2 ORCA GOAT-REACT}
GOAT-REACT is a built-in, non-ML search. A minimal input is:
\begin{verbatim}
! GOAT-REACT B3LYP D3BJ def2-SVP TightSCF
%pal nprocs 12 end
%maxcore 1100
%geom
  Constraints
    { C 0 C }
    { C 4 C }
    { C 5 C }
  end
end
* xyz 0 1
<coordinates>
*
\end{verbatim}
GOAT-REACT explores reaction pathways automatically; the reacting atoms can be held with Cartesian constraints in \texttt{\%geom}, and ORCA counts atoms from zero. For Cartesian constraints the initial position is kept and no target value is given. No \texttt{\%goat} block is used, so the defaults apply (its \texttt{MAXITER} key bounds the geometry iterations per worker, not the number of search steps). GOAT is a search tool;
its output still goes through P2--P4. It becomes expensive for large flexible
systems and is best used for channel discovery or medium-sized systems.

\section*{S5. Stop rules}
The stop rules applied at each stage are collected in Table~1 of the main text; they are reproduced here only as a pointer so that the two documents cannot drift apart.

\end{document}
